\documentclass[10pt,respuestas,a4]{examenEPN}

\usepackage{float}
\usepackage{comandosEPN}

\materia{Geometr\'ia en $\Rbb^n$}
\examen{Concurso para Técnico Docente}
%\tema{Integrales}
\autor{\textbf{Elaborado por:} Evelyn Cueva}
\fecha{20 de septiembre de 2022}
\logouno[1.6cm]{Logo_EPN.eps}
%\logodos[2.6cm]{}
\spanishdecimal{.}

\begin{document}

\encabezado

\begin{preguntas}

\item Sean $p$ y $d$ vectores en $\Rbb^3$, con $d\not=0$. Se define el conjunto 
\begin{equation}\label{eq:L}
L=\{x\in\Rbb^3\colon (x-p)\times d = 0\}.
\end{equation}
\begin{enumerate}
    \item Demuestre que $L=\{x\in\Rbb^3\colon x=p+\lambda d,\ \lambda\in\Rbb\}$, es decir, $L$ define una recta. 
    \item Sea $n\in\Rbb^3$, $n\not=0.$ ?`Qu\'e condiciones deben satisfacer $p,$ $d$ y $n$ para que  $\Pi\cap L=\emptyset,$ donde $\Pi =\{x\in\Rbb^3\colon \langle x, n\rangle = 0\}$?
\end{enumerate}

\final
 \end{preguntas}

\end{document}